%% kapa: exam=1 round=19 year=2008 period=1 form=A subject=civil_law subject_ko=민법 seq=1 number=1
%% source: Q-Net 감정평가사 기출문제 (cst00309) artlSeq=1001962 file=extracted/1008/1001962/2008년 제19회 감정평가사 1차시험 기출문제/1교시.hwp
%% conversion: kordoc 4.24.51 -> tools/gen1.py (auto); review=proofread-sonnet; flags=none
\begin{question}{1}
권리남용의 효과에 관한 설명으로 틀린 것은? (다툼이 있는 경우에는 판례에 의함)
\begin{choices}
  \item 권리의 행사가 상대방에게 고통이나 손해를 주기 위한 것이라는 주관적 요건은 권리자의 정당한 이익을 결여한 권리행사로 보여 지는 객관적인 사정에 의하여 추인할 수 있다.
  \item 채무자가 시효완성 전에 스스로 채권자의 시효중단을 현저히 곤란하게 하여 채권자가 시효중단을 할 수 없었던 경우, 채무자가 소멸시효의 완성을 주장하는 것은 권리남용으로서 허용될 수 없다.
  \item 계약효력의 주장이 권리남용에 해당하면 계약의 효력이 발생하지 않지만, 이 경우 법률관계는 위험부담의 원칙에 따라 확정된다.
  \item 권리남용이 불법행위에 해당하면 상대방은 그로 인한 손해배상을 청구할 수 있다.
  \item 항변권이 남용된 경우, 권리의 행사가 저지될 뿐 항변권 자체가 박탈되는 것은 아니다.
\end{choices}
\end{question}
\answer{3}
